% Student paper — English, typeset text, NaCl lattice as CIE clip. Compiler: XeLaTeX.
\documentclass[10pt,a4paper]{article}

\usepackage[margin=16mm,headheight=13pt,headsep=5pt,footskip=9mm]{geometry}
\usepackage{fontspec}
\usepackage[version=4]{mhchem}
\usepackage{chemfig}
\usepackage{graphicx}
\usepackage{xcolor}
\usepackage{booktabs}
\usepackage{array}
\usepackage{enumitem}
\usepackage{needspace}
\usepackage{fancyhdr}
\usepackage{amsmath,amssymb}
\usepackage{pgffor}

\raggedbottom
\setlist[enumerate]{leftmargin=1.5em,itemsep=0.12em,topsep=0.15em,parsep=0pt}
\setlist[itemize]{leftmargin=1.3em,itemsep=0.08em}

\pagestyle{fancy}
\fancyhf{}
\fancyhead[L]{\small \homeworktitle}
\fancyhead[R]{\small \homeworkmarks}
\fancyfoot[C]{\thepage}
\renewcommand{\headrulewidth}{0.3pt}

\newcommand{\homeworktitle}{Chemistry homework}
\newcommand{\homeworkmarks}{}
\newcommand{\sethomework}[2]{%
  \renewcommand{\homeworktitle}{#1}%
  \renewcommand{\homeworkmarks}{#2}%
}

\newcommand{\answerlines}[1]{%
  \par\vspace{0.08em}%
  \foreach \n in {1,...,#1}{\noindent\rule{\linewidth}{0.2pt}\\[0.62em]}%
}

\graphicspath{{images/}}
\DeclareGraphicsExtensions{.png,.jpg,.jpeg,.pdf}

\setlength{\parindent}{0pt}
\setlength{\parskip}{0.10em}
\setlength{\abovedisplayskip}{0.25em}
\setlength{\belowdisplayskip}{0.25em}

\newcommand{\qhead}[2]{%
  \par\vspace{0.28em}%
  \Needspace{4.2em}%
  \noindent{\small\textbf{#1}}\hfill{\small[#2]}\par\vspace{0.08em}%
}

% Diagram clip: stem is on the scan. Measure so label + image stay together.
\newcommand{\clipq}[3]{%
  \par
  \sbox0{\includegraphics[width=0.88\linewidth,keepaspectratio]{#3}}%
  \Needspace{\dimexpr\ht0+1.25em+2pt\relax}%
  \noindent{\small\textbf{#1}}\hfill{\small[#2]}\par\vspace{0.06em}%
  \noindent\usebox0\par
}

\newcommand{\optline}[1]{%
  \par\vspace{0.06em}{\small #1}\par
}

\newcommand{\sourcefoot}[1]{%
  {\footnotesize\color{gray}#1}\par
}

\begin{document}
\sethomework{AS Chemistry homework · 3.2 Ionic bonding}{17 marks}

\begin{center}
{\large\bfseries AS Chemistry homework · 3.2 Ionic bonding}\\[0.12em]
{\small CIE 9701 AS Chemistry · 17 marks · about 40 minutes}
\end{center}

{\small Topic: 3.2 Ionic bonding --- electrostatic attraction between oppositely charged ions, and describing ionic bonding in \ce{NaCl}, \ce{MgO} and related giant ionic compounds.
Answer in English. Section A is Paper 1 style (1 mark each). Section B is written; use CIE language (electrostatic attraction, cations/anions, giant ionic lattice, mobile ions when molten).}

\vspace{0.15em}
{\small
\begin{tabular}{@{}llr@{}}
\toprule
\textbf{Section} & \textbf{Content} & \textbf{Marks} \\
\midrule
A & Multiple choice & 7 \\
B1 & 2022 ON Paper 22 Question 1(a) & 2 \\
B2 & 2021 MJ Paper 22 Question 2(a) & 3 \\
B3 & 2025 MJ Paper 21 Question 1(a)(i) and (iii) & 2 \\
B4 & 2022 ON Paper 23 Question 3(a)(i) & 1 \\
B5 & 2025 MJ Paper 23 Question 2(c)(ii) & 2 \\
\midrule
 & \textbf{Total} & \textbf{17} \\
\bottomrule
\end{tabular}}

\vspace{0.25em}
{\large\bfseries Section A --- Multiple-choice questions (7 marks)}\par

% A1 compact
\qhead{A1.}{1}
{\small Which type of bonding is \textbf{never} found in elements?}
\par\vspace{0.08em}
{\small
\begin{tabular}{@{}p{0.46\linewidth}p{0.46\linewidth}@{}}
\textbf{A} covalent & \textbf{B} ionic \\[0.15em]
\textbf{C} metallic & \textbf{D} van der Waals' forces
\end{tabular}}
\par
\sourcefoot{2021 MJ Paper 13 Question 6}

% A2 typeset table
\Needspace{16em}
\qhead{A2.}{1}
{\small The atomic radii and ionic radii for three elements in Period 3 are shown.}

\vspace{0.12em}
\centering
{\small
\begin{tabular}{lcc}
\toprule
 & atomic radius / nm & ionic radius / nm \\
\midrule
element X & 0.118 & 0.053 \\
element Y & 0.099 & 0.180 \\
element Z & 0.160 & 0.072 \\
\bottomrule
\end{tabular}}
\par\raggedright

\vspace{0.10em}
{\small Using this data, which statement is correct?}
\begin{enumerate}[label=\textbf{\Alph*}, leftmargin=1.7em, itemsep=0.08em]
\item Element X has lower electrical conductivity than element Y.
\item Element Y has a higher melting point than element X.
\item Element Y and element Z react to form an ionic compound.
\item Element Z forms ionic compounds by gaining electrons.
\end{enumerate}
\sourcefoot{2025 ON Paper 11 Question 16}

% A3 lattice diagrams --- keep CIE clip (stem typeset; Q8 number cropped)
\Needspace{18em}
\qhead{A3.}{1}
{\small Which diagram represents the lattice structure of sodium chloride?}
\par\vspace{0.06em}
\noindent\includegraphics[width=0.88\linewidth,keepaspectratio]{img-003.png}\par
\sourcefoot{2025 MJ Paper 11 Question 8}

% A4 coordination table
\Needspace{12em}
\qhead{A4.}{1}
{\small In the sodium chloride lattice the number of chloride ions that surround each sodium ion is called the coordination number of the sodium ions.}

{\small What are the coordination numbers of the sodium ions and the chloride ions in the sodium chloride lattice?}

\vspace{0.12em}
\centering
{\small
\begin{tabular}{ccc}
\toprule
 & coordination number of sodium ions & coordination number of chloride ions \\
\midrule
\textbf{A} & 4 & 6 \\
\textbf{B} & 6 & 4 \\
\textbf{C} & 6 & 6 \\
\textbf{D} & 8 & 6 \\
\bottomrule
\end{tabular}}
\par\raggedright
\sourcefoot{2023 MJ Paper 11 Question 6}

% A5 compact formulae
\qhead{A5.}{1}
{\small Which oxide has a simple structure rather than a giant structure?}
\optline{\textbf{A} \ce{MgO}\qquad \textbf{B} \ce{Al2O3}\qquad \textbf{C} \ce{SiO2}\qquad \textbf{D} \ce{P4O10}}
\sourcefoot{2025 ON Paper 13 Question 18}

% A6 compact formulae
\qhead{A6.}{1}
{\small Which solid compound has both ionic and covalent bonding but \textbf{not} coordinate bonding?}
\optline{\textbf{A} \ce{Al2Cl6}\qquad \textbf{B} \ce{CH3COONa}\qquad \textbf{C} \ce{MgCl2}\qquad \textbf{D} \ce{NH4Cl}}
\sourcefoot{2025 MJ Paper 14 Question 12}

% A7 statements + combo key
\Needspace{16em}
\qhead{A7.}{1}
{\small The following statements compare some properties of the oxides and chlorides of sodium and silicon.}

{\small Which statements are correct?}
\begin{enumerate}[label=\arabic*., leftmargin=1.6em, itemsep=0.06em]
\item The melting point of \ce{SiCl4} is higher than the melting point of \ce{SiO2}.
\item \ce{SiCl4} reacts with water, \ce{NaCl} does not.
\item The melting point of \ce{SiCl4} is lower than the melting point of \ce{NaCl}.
\end{enumerate}

{\small The responses \textbf{A} to \textbf{D} should be selected on the basis of}

\vspace{0.10em}
\centering
{\small
\begin{tabular}{>{\centering\arraybackslash}p{0.22\linewidth}>{\centering\arraybackslash}p{0.22\linewidth}>{\centering\arraybackslash}p{0.22\linewidth}>{\centering\arraybackslash}p{0.22\linewidth}}
\toprule
\textbf{A} & \textbf{B} & \textbf{C} & \textbf{D} \\
\midrule
1, 2 and 3 are correct & 1 and 2 only are correct & 2 and 3 only are correct & 1 only is correct \\
\bottomrule
\end{tabular}}
\par\raggedright
{\footnotesize No other combination of statements is used as a correct response.}\par
\sourcefoot{2021 ON Paper 12 Question 35}

\vspace{0.35em}
{\large\bfseries Section B --- Structured questions (10 marks)}\par\vspace{0.1em}

{\normalsize\bfseries B1 · 2022 October/November Paper 22 Question 1(a) (2 marks)}\par

{\small Species such as \ce{NH4+}, \ce{CO3^2-} and \ce{PO4^3-} are examples of molecular ions.}

\qhead{B1(a)}{2}
{\small Ionic and covalent bonds both involve an electrostatic attraction between different species.}

{\small Identify the species that are electrostatically attracted to one another in:}
\begin{itemize}[leftmargin=1.2em, itemsep=0.35em]
\item an ionic bond\\[-0.15em]
\rule{\linewidth}{0.2pt}
\item a covalent bond.\\[-0.15em]
\rule{\linewidth}{0.2pt}
\end{itemize}

\vspace{0.25em}
{\normalsize\bfseries B2 · 2021 May/June Paper 22 Question 2(a) (3 marks)}\par

{\small The strength of interaction between particles determines whether the substance is a solid, liquid or gas at room temperature.}

{\small Lithium sulfide, \ce{Li2S}, is a crystalline solid with a melting point of $938\,^{\circ}$C. It conducts electricity when it is molten.}

\qhead{B2(a)(i)}{1}
{\small Give the formulae of the particles present in solid lithium sulfide.}
\answerlines{1}

\qhead{B2(a)(ii)}{2}
{\small Explain, in terms of the structure of the crystalline solid, why lithium sulfide has a high melting point.}
\answerlines{2}

\vspace{0.25em}
\Needspace{14em}
{\normalsize\bfseries B3 · 2025 May/June Paper 21 Question 1(a)(i) and (iii) (2 marks)}\par

{\small Solid sodium conducts electricity. Sodium oxide conducts electricity when molten but not when solid.}

\qhead{B3(a)(i)}{1}
{\small Name the type of bonding present in sodium and in sodium oxide.}

\vspace{0.08em}
{\small bonding in sodium \dotfill}

\vspace{0.20em}
{\small bonding in sodium oxide \dotfill}

\qhead{B3(a)(iii)}{1}
{\small Explain why sodium oxide conducts electricity when molten.}
\answerlines{2}

\vspace{0.25em}
{\normalsize\bfseries B4 · 2022 October/November Paper 23 Question 3(a)(i) (1 mark)}\par

{\small Some of the common chlorides of Period 3 elements are shown in the list.}

\vspace{0.08em}
\centering
{\small \ce{NaCl}\qquad \ce{MgCl2}\qquad \ce{AlCl3}\qquad \ce{SiCl4}\qquad \ce{PCl5}}
\par\raggedright

\qhead{B4(a)(i)}{1}
{\small From this list, identify all the chlorides that have giant ionic structures in the solid state.}
\answerlines{1}

\vspace{0.25em}
{\normalsize\bfseries B5 · 2025 May/June Paper 23 Question 2(c)(ii) (2 marks)}\par

{\small The chemical properties of oxides are related to the chemical bonding present in these compounds.}

{\small Table 2.1 shows the melting points of some oxides.}

\vspace{0.12em}
\centering
{\small
\begin{tabular}{lc}
\toprule
oxide & melting point / $^\circ$C \\
\midrule
\ce{SO2} & $-73$ \\
\ce{H2O} & 0 \\
\ce{SO3} & 17 \\
\ce{SiO2} & 1610 \\
\ce{Na2O} & 1132 \\
\ce{MgO} & 2852 \\
\ce{Al2O3} & 2072 \\
\bottomrule
\end{tabular}}
\par\vspace{0.04em}
{\small Table 2.1}\par
\raggedright

\qhead{B5(c)(ii)}{2}
{\small A student suggests the following hypothesis.}

\vspace{0.08em}
{\small\textit{The melting point of an ionically bonded oxide is only determined by the charge on the cation.}}

\vspace{0.10em}
{\small Use Table 2.1 to deduce if this hypothesis is true or false or if there is \textbf{not} enough information to make a conclusion. Explain your answer.}
\answerlines{3}

\end{document}
